\documentclass[../../../../main.tex]{subfiles}
\begin{document}

\begin{flushright}
\footnotesize\textit{【自動文字起こし・要確認】}
\end{flushright}

［解］ 平均値の定理から，$h(x) = \frac{f(x) - f(a)}{x - a}$ に対し，$h(x) = f'(c)$ なる $c$ が $a < c < x$ にある。 \dots ①

又，$f''(x) > 0$ から $f'(x)$ は単調増加。 \dots ②

$a < x < b$ のとき，
\begin{align*}
h'(x) > 0 &\iff f'(x)(x-a) - [f(x) - f(a)] > 0 \\
&\iff f'(x) > \frac{f(x) - f(a)}{x - a} = f'(c) \quad \dots \text{③}
\end{align*}
だが，②及び $c < x$ から③は明らか。よって $h'(x) > 0$ つまり $h(x)$ は単調増加。

\end{document}
